\subsubsection{Degree one}

\begin{theorem}[The causal order complex is a cochain complex]\label{thm:causal-cochain-complex}\lean{Infomaterial.Relational.Spacetime.Causal.Cohomology.causal_cochain_complex}\leanok\uses{thm:cotangent-microsupport, thm:cone-field-geometry}
The explicit cohomology the cotangent layer needed.
\cref{thm:cotangent-microsupport} introduced the degree-0 coboundary
$\delta^0$ but no $\delta^1$, no cocycle condition, no $H^1$. Here the causal
poset's nerve is given its cochain complex: $C^0$ integer functions on events,
$C^1$ on causal pairs $a \leq b$, and $C^2$ on composable triples
$a \leq b \leq c$, with simplicial coboundaries
$\delta^0 f(a,b) = f(b) - f(a)$ and
$\delta^1 \varphi(a,b,c) = \varphi(b,c) - \varphi(a,c) + \varphi(a,b)$.
The identity $\delta^1 \circ \delta^0 = 0$ telescopes, so
$\operatorname{im}\delta^0 \subseteq \ker\delta^1$ and
$H^1 = \ker\delta^1 / \operatorname{im}\delta^0$ is well-defined. This is the
causal-set realization of the abstract \v{C}ech $H^1$ of
\cref{thm:h1-vanishes}: the $\mathrm{cl\text{-}inv}$ is a 1-cocycle class,
the wavefront is its support, and an empty wavefront forces a trivial
coboundary class.
\end{theorem}
